\documentclass[11pt]{article}
\usepackage[utf8]{inputenc}
\usepackage[T1]{fontenc}
\usepackage{geometry}
\usepackage{enumitem}
\usepackage{amsmath}
\usepackage{xcolor}

\geometry{margin=1in}

\title{Week 5 Quiz: Controlling Style, Tone, and Format}
\author{Prompt Engineering Course}
\date{}

\begin{document}

\maketitle

\noindent \textbf{Name:} \underline{\hspace{8cm}} \hspace{1cm} \textbf{Date:} \underline{\hspace{3cm}}

\vspace{0.5cm}
\noindent \textbf{Instructions:} Answer all questions by circling the best answer for each multiple choice question.

\vspace{0.5cm}

\begin{enumerate}

\item What are the three dimensions of linguistic register?
\begin{enumerate}[label=\Alph*.]
    \item Formal, casual, neutral
    \item Field, tenor, mode
    \item Syntax, lexicon, rhythm
    \item Voice, tone, style
\end{enumerate}

\item Which of the following best describes ``tenor'' in linguistic register?
\begin{enumerate}[label=\Alph*.]
    \item The technical complexity of vocabulary used
    \item The relationship and power distance between communicators
    \item Whether communication is spoken or written
    \item The average sentence length in a document
\end{enumerate}

\item In Brand Voice Fingerprinting, what does ``syntax signature'' measure?
\begin{enumerate}[label=\Alph*.]
    \item Preferred vocabulary choices
    \item Sentence length, clause complexity, and passive voice percentage
    \item The rhythm and pacing of writing
    \item The point of view used in writing
\end{enumerate}

\item Why should you provide both a strong sample AND a counter-example when doing Voice and Tone Diagnostics?
\begin{enumerate}[label=\Alph*.]
    \item To confuse the AI model
    \item AI learns boundaries better from contrasts than from positive examples alone
    \item Counter-examples are required by law
    \item To make the prompt longer and more detailed
\end{enumerate}

\item What is a ``structural contract'' in prompt engineering?
\begin{enumerate}[label=\Alph*.]
    \item A legal agreement between you and the AI provider
    \item A specification that defines exactly how output should be formatted for downstream processing
    \item A contract stating how many prompts you can use
    \item A formal tone requirement
\end{enumerate}

\item Which prompt instruction would create a more ``urgent'' tone?
\begin{enumerate}[label=\Alph*.]
    \item ``Use passive voice and hedge words''
    \item ``Write in third person with long, complex sentences''
    \item ``Use short sentences, action verbs, and explicit deadlines''
    \item ``Avoid all time references''
\end{enumerate}

\item In multi-stakeholder voice reconciliation, what is the recommended approach when different departments have conflicting voice requirements?
\begin{enumerate}[label=\Alph*.]
    \item Choose the CEO's preference
    \item Use layered prompts with priority matrices
    \item Ignore the conflicts
    \item Average all requirements together
\end{enumerate}

\item What is the purpose of programmatic tone testing?
\begin{enumerate}[label=\Alph*.]
    \item To manually review every AI output
    \item To automatically validate tone consistency using metrics like sentiment analysis and readability scores
    \item To test if the AI is connected to the internet
    \item To count the number of words in output
\end{enumerate}

\item Which statement about multi-dimensional register specifications is TRUE?
\begin{enumerate}[label=\Alph*.]
    \item ``Be professional'' is just as precise as ``CEFR B1 + peer tenor + written-planned mode''
    \item Multi-dimensional specs are too complicated to be useful
    \item Multi-dimensional specs provide precise, reproducible style control
    \item Vague instructions always work better than detailed specifications
\end{enumerate}

\item What does ``mode'' control in linguistic register?
\begin{enumerate}[label=\Alph*.]
    \item The relationship between communicators
    \item Technical vocabulary density
    \item Whether communication is spoken/written and planned/spontaneous
    \item The emotional tone of the message
\end{enumerate}

\end{enumerate}

\end{document}
